\section{一般 $\t$ 处的分析}\label{analog}

\begin{keytakeaways}
\begin{itemize}\setlength{\itemsep}{2pt}
  \item \textbf{物理图像}：把四维方程形变到一般的 $t$ 后，$y\ne0$ 处的全纯数据不再由希格斯丛给出，而由复平坦联络及其上的 oper 结构给出；纽结（'t Hooft 算符）则对应 oper 的奇点。
  \item \textbf{数学工具}：Hitchin 方程的解模空间作为 hyper-Kähler 流形、其由 $\CP^1$ 上参数 $\zeta$ 参数化的复结构族、复平坦联络，以及 oper（平坦 $SL(2,\C)$ 丛连同处处不被 $D_z$ 平行移动保持的全纯线子丛）。
  \item \textbf{主要结论}：$t=1$ 的旧图景对应例外值 $\zeta=0$ 的希格斯丛极限；一般 $\zeta\ne0$ 时三个对易算符恰描述 $C\times\R_+$ 上的复平坦联络；Nahm 极边界条件在 $y=0$ 附近精确对应二维 oper，而带指数 $k$ 的奇点对应磁荷 $k$ 的 't Hooft 算符。
  \item \textbf{遗留问题}：矩映射条件的可解性只是"猜测性地"成立；奇点处的模型解目前只能数值求得（见附录）。
\end{itemize}
\end{keytakeaways}

\subsection{一些基本事实}\label{motivation}

我们仍在为任意磁荷的 't Hooft 算符寻找有用的描述，因此考虑把我们试图求解的方程形变到另一个方程，后者在沿某个方向拉伸时可能表现得更方便。

存在有用形变的一个强有力提示，来自把推广 Bogomolny 方程 (\ref{eqform})、(\ref{mmtmap}) 与 Hitchin 方程 \cite{Hitchin} 联系起来。Hitchin 方程是有向二维流形 $C$ 上关于联络 $A$ 与伴随表示值一形式 $\phi$ 的方程：
\begin{align}\label{hitcheq} F-\phi\wedge\phi & = 0 \cr \d_A\star\phi& = 0 \cr
\d_A\phi&=0.\end{align}
这些方程的解模去规范变换后的模空间，是一个 hyper-Kähler 流形 $\M_H(G,C)$。作为一个 hyper-Kähler 流形，$\M_H(G,C)$ 有一族由 $\Bbb{CP}^1$ 的一个拷贝参数化的复结构。$\M_H(G,C)$ 的复结构有一个简单的解释。它们对应于把 (\ref{hitcheq}) 中的三个实方程分成两部分的方式：两个实方程（可以合并为一个复方程），以及一个与前两者“正交”的第三个实方程。复方程于是描述在其中一个复结构下参数化 $\M_H(G,C)$ 的全纯数据，而剩下的实方程是一个矩映射条件。

一个例外分解对应于这样的情形：复方程由 (\ref{hitcheq}) 中最后两个方程合并而成，从而得到希格斯丛方程
\begin{equation}\label{higgs}\bar\partial_A\varphi=0.\end{equation}
更一般的分解涉及一个复参数 $\zeta$。对任何 $\zeta$，令
\begin{align}\label{generic} \D_z^\zeta& = \frac{D}{D z}-\zeta^{-1}[\phi_z,\,\cdot\,]\cr
                \D_{\bar z}^\zeta &=\frac{D}{D\bar z}+\zeta[\phi_{\bar z},\,\cdot\,].\end{align}
（我们在这里用 $\phi_z$ 与 $\phi_{\bar z}$ 代替 $\varphi$ 与 $\bar\varphi$。）对任何 $\zeta$，方程
\begin{equation}\label{nurco}[\D_z^\zeta,\D_{\bar z}^\zeta]=0 \end{equation}
等价于 Hitchin 方程 (\ref{hitcheq}) 的两个实线性组合。$\M_H(G,C)$ 上可能的复结构由 $\zeta$ 参数化，其中我们在 $\zeta$ 平面上加一个无穷远点使之成为 $\Bbb{CP}^1$。对每个 $\zeta$，定义一个复结构 $I_\zeta$，在其中方程 (\ref{nurco}) 被视为支配全纯数据的方程；在此复结构中，全纯变量为 $\A^\zeta_z=A_z-\zeta^{-1}\phi_z$ 与 $\A_{\bar z}^\zeta=A_{\bar z}+\zeta \phi$，而方程 (\ref{nurco}) 关于这些变量是全纯的。方程的第三个线性组合在复结构 $I_\zeta$ 中被视为矩映射条件。
% [Tutor Note: 易错点] zeta 的参数化约定与不少文献相反（常见写法是 D_z 中带 +zeta 而 D_zbar 中带 -zeta^{-1}）；若照搬别的文献，复结构 I_zeta 会整体换成 zeta→1/zeta。文中特别指出 zeta=0 与 zeta=∞ 是例外值：此时不能再把方程写成"复联络平坦"这一条件，而分别退化为希格斯丛方程及其复共轭。

对一般的 $\zeta$，方程 (\ref{generic}) 简单地说就是复联络 $\A^\zeta=\A_z^\zeta \d z+\A_{\bar z}^\zeta \d \bar z$ 是平坦的。一旦方程 (\ref{generic}) 用 $\A^\zeta$ 表达，它就不显含 $\zeta$。因此，当 $\M_H(G,C)$ 被简单地视为复结构 $I_\zeta$ 下的复流形、而不考虑其 Kähler 度规时，它对一般的 $\zeta$ 与 $\zeta$ 无关。$\zeta$ 的例外值是 0 与 $\infty$。例如，对 $\zeta\to 0$，我们必须把 $\D_z^\zeta$ 乘以 $\zeta$，于是它约化为 $-[\phi_z,\,\cdot\,]$。与此同时，对 $\zeta=0$，$\D_{\bar z}=D_{\bar z}$。所以方程 (\ref{nurco}) 的 $\zeta\to 0$ 极限就是希格斯丛方程 (\ref{higgs})。由类似的推理，(\ref{nurco}) 的 $\zeta\to\infty$ 极限是 (\ref{higgs}) 的复共轭。
% [Tutor Note: 易错点] zeta→0 的极限须先把 D_z^zeta 乘以 zeta 再取极限，否则其中 zeta^{-1} phi_z 项发散；zeta→∞ 时则须先把 D_zbar^zeta 乘以相关的 zeta 因子。这与第一部分中 t→0、t→∞ 处要把第二个方程乘以 t 或把第一个方程乘以 t^{-1} 的处理完全同构，两处缩放因子极易混淆，核对时请以各自的缩放为准。

希格斯丛方程在我们的分析中起了突出的作用，因为第~\ref{preliminaries}~节中我们从方程 $[\D_1,\D_2]=0$ 导出了 $y\not=0$ 处的全纯数据，而后者不是别的，正是希格斯丛方程。（那一分析中的矩映射条件并不是希格斯丛通常的矩映射方程——即 (\ref{hitcheq}) 中的第一个方程——而是三维方程 (\ref{mmtmap})。）至此，我们不禁要问：能否修改导出 (\ref{eqform}) 的构造，使得 $y\not=0$ 处的全纯数据不再由希格斯丛给出，而由 Hitchin 方程另一个复线性组合的解给出？事实上这可以做到。不过，做到这一点的途径并非唯一确定。

有一个特别好的参数可以用来改变底层的四维方程，它简单地就是式 (\ref{bpseqns}) 中的参数 $\t$。除了改变方程之外，我们还应当考虑改变边界条件的可能性。在六个场 $(\vec A,\vec\phi)$（即 $A$ 与 $\phi$ 中与边界相切的分量）中，只有 $\vec\phi$ 在基本 Nahm 极解 (\ref{nahmpole}) 中于 $y=0$ 处有奇点。然而，一旦丢掉对空间坐标 $\vec x$ 的依赖，$\N=4$ 超 Yang-Mills 理论就具有一个旋转 $\vec A$ 与 $\vec \phi$ 六个分量的 $SO(6)$ 对称性。因此，人们可以用对 Nahm 极施加一个 $SO(6)$ 旋转所得的解来满足经典 Yang-Mills 方程（这样的旋转 Nahm 极曾在 \cite{gw} 第 4 节中研究过）。当然，这一旋转一般会改变未破缺的超对称。然而事实证明（详见附录 \ref{fromsix} 与 \ref{foureight}），只要旋转矩阵包含在某个 $SU(3)$ 子群（属于 $SO(6)$）中，未破缺的超对称就不变；而若它包含在某个 $U(3)$ 子群中，则未破缺的超对称会以对应于四维方程 (\ref{bpseqns}) 中参数 $\t$ 改变的方式改变。所以在旋转边界条件方面有相当大的自由度，无论是否改变 $\t$。
% [Tutor Note: 约定差异] 用 SO(6) R 对称性旋转 Nahm 极时，落在 SU(3) 子群内的旋转保持超对称不变，而落在 U(3)（但不含于 SU(3)）内的旋转等价于改变 t：这是因为 U(3)/SU(3)=U(1) 正对应 t 所参数化的那个自由度。t 与旋转角度的精确关系见附录（标签 fromsix、foureight），正文只给出定性结论。

在按上述方式作出选择之后，形变后的方程与边界条件有一个三维约化，其形式就是熟悉的 (\ref{commuting})、(\ref{momentmap})，但 $\D_i$ 的定义不同。把大部分细节留到附录，我们先总结在某个有启发性的特殊情形下出现的一些公式。若希望保持边界的 $SO(3)$ 旋转对称性，则 Nahm 极的旋转只能依赖于一个参数：$\vec A$ 的极向部分必须等于常数 $\zeta$ 乘以 $\vec\phi$ 的极向部分。事实证明，这样的边界条件与四维方程 (\ref{bpseqns}) 相容，只是 $\t$ 的值被修改为依赖于 $\zeta$。

在约化到三维时，我们现在必须施加 $A_1-\zeta\phi_1=0=\phi_y$。三个对易的微分算符 $\D_i$ 于是具有形式
\begin{align}\label{neqformo}\D_1 & =2\frac{D}{D\bar z}  +2 \zeta [\bar \varphi,\,\cdot\,] \cr
                                                              \D_2 & = - 2 \zeta \frac{D}{D z} + 2[\varphi,\,\cdot\,]\cr
                                                               \D_3& = \frac{\D}{\D y}. \end{align}
这里，如附录所述，$\A_y$ 是 $A_y$ 与 $\phi_1$ 的适当线性组合。$\D_i$ 对易的条件必须补充以一个矩映射条件，后者也在附录中描述。正如保持四个超荷的几何所预期的那样，对易性约束 $[\D_i,\D_j]=0$ 可以由一个超势导出，而该超势实际上是 Chern-Simons 泛函的一个倍数：
\begin{equation}\label{supertwo}\W=\frac{1}{4\pi i}\int_{W}\Tr\,\left(\A\wedge\d \A+\frac{2}{3}\A\wedge
\A\wedge\A\right).\end{equation}
更准确地说，在变分 $\W$ 以导出方程 $[\D_i,\D_j]=0$ 时，要对 $\A$ 在 $y=0$ 处的变分施加一个约束，以避免 $\W$ 的变分中出现 delta 函数项。

远离 $y=0$ 处的全纯数据现在容易描述。若 $\zeta=0$，我们恢复第~\ref{teqone}~节所研究的图景；但一旦 $\zeta$ 非零（且不为无穷），三个对易算符 $\D_i$ 就简单地描述 $C\times \R_+$ 上的一个复平坦联络。该平坦联络的协变导数为 $\D_z=\D_1/2$、$\D_{\bar z}=-\D_2/2\zeta$、$\D_y=\D_3$。所以对易算符 $\D_i$ 就简单地描述一个复平坦联络。我们将把这个平坦联络称为 $\A$，而不管它原本如何用 $A$ 与 $\phi$ 定义。$\A$ 受一个矩映射条件约束，该条件与最常研究的复平坦联络的矩映射条件 \cite{Corlette} 并不完全一致，尽管定性上相似；我们期望它对任何 $\zeta$ 都有唯一解。

因此，远离 $y=0$ 处的全纯数据就是一个 $C\times\R_+$ 上的复平坦联络，或者等价地（由于此空间可缩到 $C$），是 $C$ 上的一个复平坦联络。我们主要关心 $C$ 为 $\R^2$ 或 $\Bbb{CP}^1$ 的情形。在这两种情形下，$C$ 都是单连通的，所以 $C$ 上的复平坦联络是平凡的。人们因此可能疑惑：还能发生什么有趣的事情？答案是：大部分有趣的结构将来自 $y=0$ 处的边界条件。在对称性破缺的情形下，$y=\infty$ 处的行为中也有一些有趣的结构。

\subsection{Nahm 极与 oper}\label{flatpole}

第一点是理解一个复平坦联络如何能具有 Nahm 极。答案是：极点出现在 $\A_y$ 与 $\A_z$ 中，但不出现在 $\A_{\bar z}$ 中。带 Nahm 极的平坦联络的模型例子是
\begin{align}\label{modsol}\A_z&=\frac{\ct_+}{y}\cr
\A_{\bar z} & = 0 \cr
\A_y & = \frac{\ct_3}{y}.   \end{align}
只要 $[\ct_3,\ct_+]=\ct_+$，这就描述一个平坦联络。我们感兴趣的是 $\ct_+=\ct_1+i\ct_2$ 的情形，其中 $\ct_i$（$i=1,2,3$）生成主 $\frak{su}(2)$ 子代数（$\frak g$ 的）。

作为平坦联络，它可以用公式 $\d+\A=g\d g^{-1}$ 描述。例如对 $\frak{su}(2)$，我们可以显式地取
\begin{equation}\label{inggauge}
g =\begin{pmatrix} y^{-1/2} & -zy^{-1/2} \cr 0 & y^{1/2} \end{pmatrix},
\end{equation}
这给出
\begin{align}\label{odsol} \A_z& = \begin{pmatrix}0&1\cr 0 & 0 \end{pmatrix}\frac{1}{y}\cr
                            \A_{\bar z}&=0 \cr
                            \A_y& = \begin{pmatrix}1 & 0 \cr 0 & -1\end{pmatrix}\frac{1}{2y}.\end{align}
或者，模型解也可以由非奇异的平坦联络
\begin{equation}\label{nodsol}\A_z=\begin{pmatrix}0 & 1\cr 0 & 0\end{pmatrix}, ~\A_{\bar z}=\A_y=0
\end{equation}
经与 (\ref{singgauge}) 相同的奇异规范变换生成：
\begin{equation}\label{winggauge}
g =\begin{pmatrix} y^{-1/2} & 0 \cr 0 & y^{1/2} \end{pmatrix}.
\end{equation}

现在我们的做法与希格斯丛情形颇为相似，以解释必须对复平坦丛 $E$ 施加什么条件，才能使它（模去奇异性更低的项）在 $y=0$ 附近化为形式 (\ref{odsol})。对某个固定的 $y>0$，我们把 $E$ 到 $C=C\times y$ 的限制简单地记作 $E$。
 我们考虑方程 $\D_ys=0$ 的、在 $s\to 0$ 时按 $y^{1/2}$ 消失的解 $s$。满足这些条件的截面 $s$ 张成一个秩一子丛 $\L\subset E$。在模型解的情形下，任何这样的 $s$ 都是
\begin{equation}s=y^{1/2}\begin{pmatrix}0 \cr 1\end{pmatrix},\end{equation}
的倍数，因此如希格斯丛情形一样，$\L$ 简单地由上分量为零的截面张成。
% [Tutor Note: 原文笔误] 本段有两处笔误：一是"把 E 到 C=C×y 的限制记作 E"，按上下文应为限制到 C×{y}（可对照第一部分的同一句话）；二是"在 s→0 时按 y^{1/2} 消失"，应为在 y→0 时按 y^{1/2} 消失，否则条件没有意义。译文按原文保留了这两处公式写法（C=C×y 与 s→0）未作改动。

同样如前面那样，若把 $E$ 看作 $C$ 上的平坦丛，则 $\L$ 是一个全纯子丛；的确，我们刚刚定义的对象 $s$ 是 $\L$ 的全纯截面，因为它当然被 $\D_{\bar z}$ 湮灭。然而，$s$ 被 $\D_z$ 湮灭并不成立。相反，看一眼前面的公式就立刻看出
\begin{equation}\label{tomox}  s\wedge \D_z s = 1.    \end{equation}

这就把我们引向“oper”这一数学概念。（关于这一概念的解释，以及对我们后文中将遇到的许多相关想法的综述，见 \cite{Frenkel} 或 \cite{Teschner}。）对 $G=SU(2)$，oper 是黎曼面 $C$ 上一个结构群为 $SL(2,\C)$ 的秩二平坦复丛 $E$，连同具有下述性质的全纯线子丛 $\L\subset E$：$\L$ 处处不被 $\D_z$ 的平行移动所保持。最后这句话的意思是：若 $s$ 是 $\L$ 的一个局部非零全纯截面，则 $\D_zs$（由于 $\D_z$ 是 $E$ 上的联络，它取值于 $E$）处处不取值于 $\L$。等价的陈述是（由于 $\L$ 由 $s$ 的倍数张成）：$\D_z s$ 处处不是 $s$ 的倍数。或者等价地，$s\wedge \D_z s$ 处处非零。最后这一陈述不依赖于非零截面 $s$ 的选取，因为若我们把 $s$ 换成 $f s$（其中 $f$ 是 $C$ 上的非零局部全纯函数），就有 $s\wedge \D_zs \to f^2 s\wedge \D_zs$。
% [Tutor Note: 物理图像] oper 条件的几何含义是：全纯线子丛 L 与复结构相容，但沿 D_z 方向被"扭转"到最大程度——处处不与自身平行。用希格斯丛的语言，这正对应 s∧φs 处处非零（零点阶数为零）的情形；而当 s∧D_z s 在某些点按 (z−z_a)^{k_a} 消失时，就得到"带奇点的 oper"，其奇点恰好编码边界上的 't Hooft 算符。换言之，第一部分中"希格斯丛 + 线子丛"的图景在一般 t 下被"复平坦丛 + oper"取代。

我们已经从 Nahm 极边界条件中提取出一个 oper 结构；反过来，让我们看看：给定 $C$ 上的一个 oper，即满足刚刚描述的条件的一对 $(E,\L)$，我们能够构造出 Nahm 极边界条件的一个解。取一个使 $\L$ 由上分量为零的向量张成的规范。$\L$ 的全纯性意味着在此规范下 $\A_{\bar z}$ 是下三角的。对任何非零的局部截面 $s$，oper 条件 $s\wedge\D_z s\not=0$ 意味着在此规范下 $\A_z$ 的右上角矩阵元非零。通过进一步的对角规范变换，我们可以把这个矩阵元置为 1。然后我们把带联络 $\A$ 的平坦丛 $E$ 从 $C$ 拉回到 $C\times \R_+$（以得到 $C\times\R_+$ 上一个不依赖于 $y$ 的平坦联络），并作奇异规范变换 (\ref{winggauge})。右上角矩阵元为 1 意味着在奇异规范变换之后，$\A_z$ 具有模型解 (\ref{odsol}) 的奇异行为；由于是下三角的，$\A_{\bar z}$ 不获得奇性。最后，规范变换恰好使 $\A_y$ 具有模型解的形式。

所以我们至少在 $G$ 秩为 1 的情形下证明了：二维 oper 恰好对应于三维中模去奇异性更低的项后的 Nahm 极边界条件的解。猜测性地，通过进一步的、光滑的复值规范变换，可以满足矩映射条件。
% [Tutor Note: 含糊点拨] "猜测性地（conjecturally）通过进一步光滑复规范变换满足矩映射条件"是关键但未证明的一步：其证明需要 oper 所对应的希格斯丛满足稳定性，并援引 Hitchin 的矩映射可解性定理（与第一部分"我们期望可以解出"处是同一类论证）。作者用 conjecturally 而非 "we expect"，语气更弱，阅读时不可当作结论。

\subsubsection{若干进一步的注记}\label{somefurth}
我们补充以下技术性注记。由于我们希望考虑最小磁荷的 't Hooft 算符，在秩为 1 的情形下我们将取规范群为 $G=SO(3)$，而不是 $SU(2)$。相应地，我们应当用伴随丛 $\ad(E)$ 而不是 $E$ 来重新表述上面的推导。由于我们的考虑在 $C$ 上是局域的，用伴随丛重写这一构造不会有太大变化，我们将略过。（主要差别在于：一般能整体定义的并非 $\L$ 而是 $\L^2$，它是 $\ad(E)$ 的一个子丛。）

此外，我们对秩为 1 的 $G$ 所说的一切对任何半单 $G$ 都有类似物，其方式与我们在第~\ref{bc}~节末尾在希格斯丛情形下所指出的有些相似。例如，对 $G=SU(n)$，oper 是秩为 $n$、具有 $SL(n,\C)$ 值和乐的平坦复丛 $E\to C$，连同全纯线子丛 $\L\subset E$，其性质是：若 $s$ 是 $\L$ 的局部非零全纯截面，则 $s,\,\D_z s,\dots,\D_z^{n-1}s$ 给出 $E$ 的一个局部平凡化。方程 $[\D_i,\D_j]=0$ 连同 Nahm 极边界条件，通过与我们已给出的论证类似的论证，确定了这样一个 oper 结构。

\subsection{带奇点的 oper}\label{opsing}

现在我们想修改 Nahm 极边界条件，以在 $y=0$ 的边界上纳入额外的奇点——我们将把它们与 't Hooft 算符联系起来。

我们所需要的奇点类型，可以通过与希格斯丛的讨论类比而猜出。我们仍将有一个平坦的 $G_\C$ 丛 $E\to C\times \R_+$，而在一般的边界点附近，平坦联络将看起来像模型解 (\ref{modsol})，在相差一个酉（即取值于 $G$ 而非 $G_\C$ 的）规范变换的意义下。我们仍可通过考虑满足 $\D_ys=0$ 且在 $y\to 0$ 时消失的截面 $s$ 来定义一个全纯线子丛 $\L\subset E$；我们仍要求：若 $s$ 是 $s$ 的一个局部全纯截面，则 $s\wedge \D_z s$ 一般地非零。
% [Tutor Note: 原文笔误] "若 s 是 s 的一个局部全纯截面"中第二个 s 应为 L（线子丛）：对照上一节 oper 的定义"若 s 是 L 的局部非零全纯截面"，此处显然是笔误。译文按原文保留了公式中的 s 未改，读者应自行替换为 L。

唯一的差别是，现在我们假设存在例外点 $z_a,~a=1,\dots,\ss$，在那里 $s\wedge \D_z s$ 为零。事实上，我们指定正整数 $k_a $，并要求 $s\wedge \D_z s$ 在 $z\to z_a$ 时按 $k_a $ 阶消失：
\begin{equation}\label{badpoint} s\wedge \D_z s\sim (z-z_a)^{k_a}.\end{equation}
这类似于在希格斯丛情形下要求 $s\wedge\varphi s\sim (z-z_a)^{k_a}$。

要得到具有这类边界行为的经典或量子规范理论的精确问题，剩下的工作就是精确规定矩映射条件 $\sum_i[\D_i,\D_i^\dagger]=0$ 的解（或可维数约化到它的四维、五维方程的解）在 $z=z_a$ 处应当具有何种奇性。正如大多数这类问题那样，为此人们先对只有一个奇点（例如位于 $z=0$（及 $y=0$）、且 $k$ 任意）的情形求出模型解。然后要求解在每个点 $z=z_a$、$y=0$ 附近的奇异行为，当 $k=k_a$ 时，应与该模型解的行为一致。对希格斯丛的情形，相应的模型解（对秩为 1 的 $G$）已在 \cite{fiveknots} 第 3.6 节中找到；但对此处所考虑的推广，目前我们只能数值地求出模型解。它们描述于附录 \ref{bco} 中。

我们迄今描述的对象，对应于 $C\times \R_+$ 上平坦性与矩映射条件的解，其边界条件与 Nahm 极或 oper 相联系，只是在有限多个边界点上 oper 条件被修正。特别地，一旦离开 $y=0$，就简单地有一个平坦丛（带一个矩映射条件）。因此平坦丛绕点 $z=z_a$ 的单值性是平凡的，于是点 $z=z_a$ 处的例外行为只影响对 $(E,\L)$ 的 oper 性质，而不影响 $E$ 的平坦性。这类奇点被称为具有平凡单值性的 oper 奇点。
